% !TEX root = ../metatime_monograph.tex
\chapter{Electron Action}
\label{ch:electron_action}

The electron is the lightest charged lepton, with mass \(m_e =
0.510999\)~MeV.  Within the Metatime framework, this corresponds to the
largest action on the Poincar\'{e} disk:

\begin{equation}
S_e = \frac12 - 3\kappa + \frac{\kappa}{L_3} - \frac{\kappa^2}{2}.
\label{eq:Se_full}
\end{equation}

We now derive each term.

\section{Spin-1/2 Degeneracy}

The leading term \(1/2\) encodes the spin-1/2 degeneracy of the electron.
On the Poincar\'{e} disk, a spin-1/2 fermion has two accessible spin states,
contributing an additive constant to the action:

\begin{equation}
S_{\text{spin}} = \frac12.
\label{eq:spin_term}
\end{equation}

\section{Three-Generation Information Deficit}

The term \(-3\kappa\) represents the shared information deficit across the
three fermion generations.  Each generation contributes \(-\kappa\) to the
action, reflecting the information-theoretic cost of generational replication:

\begin{equation}
S_{\text{gen}} = -3\kappa.
\label{eq:gen_term}
\end{equation}

The negative sign indicates that more generations reduce the action, making
the corresponding particles heavier.

\section{Collatz Channel Correction}

The term \(+\kappa/L_3\) is a correction arising from the Collatz structure
of flavour space.  The factor \(1/L_3 = 1/7\) reflects the structural depth
of the first generation:

\begin{equation}
S_{\text{collatz}} = \frac{\kappa}{L_3}.
\label{eq:collatz_term}
\end{equation}

\section{Second-Order Curvature Correction}

The term \(-\kappa^2/2\) is the leading curvature correction on the
Poincar\'{e} disk.  It is quadratic in \(\kappa\) and represents the
information curvature of the disk:

\begin{equation}
S_{\text{curv}} = -\frac{\kappa^2}{2}.
\label{eq:curv_term}
\end{equation}

\section{Full Evaluation}

Summing all terms:

\[
\begin{aligned}
S_e &= \frac12 - 3\kappa + \frac{\kappa}{L_3} - \frac{\kappa^2}{2} \\[4pt]
    &= 0.5 - 3(0.00919315000636) + \frac{0.00919315000636}{7}
       - \frac{(0.00919315000636)^2}{2} \\[4pt]
    &= 0.5 - 0.02757945001908 + 0.00131330714377
       - 0.00004225750274 \\[4pt]
    &= 0.4736916001 \ldots
\end{aligned}
\]

Thus the electron action is:

\[
\boxed{S_e = 0.4736916001\ldots}
\]

\section{Mass Prediction}

From the exponential mass formula \(m = E_P e^{-S/\kappa}\):

\[
\begin{aligned}
m_e &= E_P \, e^{-S_e/\kappa} \\
    &= 1.22089 \times 10^{22} \cdot
       e^{-0.4736916001 / 0.00919315000636} \\
    &= 1.22089 \times 10^{22} \cdot
       e^{-51.525803\ldots} \\[4pt]
    &= 0.511642\ \mev.
\end{aligned}
\]

Comparing with the PDG value:

\[
m_e^{\text{calc}} = 0.511642\ \mev,\qquad
m_e^{\text{PDG}} = 0.510999\ \mev,\qquad
\text{error} = +0.126\%.
\]

The agreement is within 0.13\%, confirming that the action construction
captures the electron mass scale correctly.
